% !PRINTED 496
\weberpage{496}

\gesperrt{4. Die Permutationen von $m$ Ziffern zerfallen in zwei Arten, von denen die erste in eine gerade, die zweite in eine ungerade Anzahl von Transpositionen zerlegbar ist.}

Jede dieser beiden Arten umfasst gleich viel Permutationen, nämlich $\frac{1}{2}\Pi(m)$; denn durch Hinzufügung einer festen Transposition geht jede Permutation der ersten Art in eine der zweiten Art über und umgekehrt.

Die Zusammensetzung zweier Permutationen von gleicher Art giebt stets eine Permutation der ersten Art, während eine Permutation erster und zweiter Art zusammengesetzt, eine Permutation zweiter Art ergeben, wie aus der Zerlegung in Transpositionen sofort zu ersehen ist. Daraus folgt:

\gesperrt{5. Die Permutationen der ersten Art bilden eine Gruppe.}

Nennen wir eine Function von $m$ Veränderlichen $u_0, u_1 \dots u_{m-1}$ oder auch von $m$ von einander verschiedenen Grössen $\alpha_0, \alpha_1 \dots \alpha_{m-1}$, die ihr Zeichen ändert, wenn zwei der Variablen mit einander vertauscht werden, wie z.\,B. das Product aller Differenzen:
\[
(u_0 - u_1)(u_0 - u_2) \dots (u_0 - u_{m-1})
\]
\[
(u_1 - u_2) \dots (u_1 - u_{m-1})
\]
\[
\cdots\cdots\cdots\cdots
\]
eine \gesperrt{alternirende Function}, so wird durch die Permutationen der ersten Art eine solche Function nicht verändert, während sie durch eine Permutation der zweiten Art geändert, nämlich in den entgegengesetzten Werth verwandelt wird. Die alternirenden Functionen \gesperrt{gehören} also zu der Gruppe der Permutationen der ersten Art.

Wie wir früher die Gruppe, zu der die symmetrischen Functionen gehören, symmetrische Gruppe genannt haben, so nennen wir die Gruppe der Permutationen der ersten Art, zu der die alternirenden Functionen gehören, \gesperrt{alternirende Gruppe}. Die Gruppe, die aus der einzigen identischen Permutation besteht, nennen wir die \gesperrt{identische Gruppe} oder auch die \gesperrt{Einheitsgruppe}.

Eine cyklische Permutation von $n$ Gliedern lässt sich in $n - 1$ Transpositionen zerlegen, wie man aus der Zusammensetzung
\[
(1, 2, 3 \dots n) = (1, 2)(1, 3) \dots (1, n)
\]
%CONT>>
